Weather files for a particular past year are built from station records, which almost always have gaps. Filling methods are usually judged against the weather, not by their effect on simulations. We hid known hours in 17 complete Swedish weather-years and refilled them with seven methods, from linear interpolation to a nearby station and the ERA5 reanalysis, including a stepwise method with and without a switch to ERA5 for long gaps. Gaps of 3 hours to one week fell at random or during heatwaves and cold spells, and six typical buildings were simulated. The results show that gap filling barely changes annual heating: with four 12-hour gaps a year, every method stayed within 0.3\% in 95\% of cases, and the extra heating followed the building's heat-loss coefficient times the temperature error summed over the gap. Peaks are more fragile. A filled extreme is usually milder than the real one, and the year's peak often has a near-equal day, so when a week-long gap hid the heating peak, the filled peak moved by over a day in half the cases. Which weather was hidden mattered more than the method; methods differed mainly in worst cases. These came from nearby stations that matched well in ordinary weather but not during a week-long extreme; switching to ERA5 after one or two days (three hours for solar radiation) avoided them. The tested gaps cover 99\% of the gaps in usable Swedish files. Filled hours should be reported, and peaks given with an uncertainty band.
